\documentclass[aspectratio=169]{beamer}

% Tema UFS --- Universidade Federal de Sergipe
\usetheme{UFS}

% Pacotes base
\usepackage[utf8]{inputenc}
\usepackage[portuguese]{babel}
\usepackage{amsmath,amssymb}
\usepackage{graphicx}
\usepackage{booktabs}
\usepackage{xcolor}

% TikZ e bibliotecas
\usepackage{tikz}
\usetikzlibrary{shapes.geometric, arrows.meta, positioning, matrix, fit,
  backgrounds, decorations.pathreplacing, shadows, patterns, calc, trees}

% Listagens --- paleta VS Code Light+
\usepackage{listings}
\definecolor{bgcolor}{HTML}{FFFFFF}
\definecolor{linecolor}{HTML}{DDDDDD}
\definecolor{numcolor}{HTML}{999999}
\definecolor{kwcolor}{HTML}{0000FF}
\definecolor{strcolor}{HTML}{A31515}
\definecolor{cmtcolor}{HTML}{008000}

\lstdefinestyle{ufsStyle}{
  backgroundcolor=\color{bgcolor},
  basicstyle=\ttfamily\footnotesize\color{black},
  keywordstyle=\color{kwcolor}\bfseries,
  stringstyle=\color{strcolor},
  commentstyle=\color{cmtcolor}\itshape,
  numberstyle=\tiny\color{numcolor},
  numbers=left, numbersep=12pt,
  xleftmargin=20pt, framexleftmargin=20pt,
  breaklines=true, tabsize=4,
  keepspaces=true, showspaces=false,
  showstringspaces=false, showtabs=false,
  frame=single, rulecolor=\color{linecolor},
  framesep=4pt, captionpos=b, upquote=true,
  columns=fullflexible,
}
\lstset{style=ufsStyle, language=C}

\lstdefinestyle{tinyC}{
  style=ufsStyle, language=C, numbers=none,
  basicstyle=\ttfamily\scriptsize\color{black},
}
\lstdefinestyle{codC}{
  style=tinyC, basicstyle=\ttfamily\fontsize{7.5}{8.6}\selectfont\color{black},
}
\lstdefinestyle{shellStyle}{
  style=ufsStyle, language=bash, numbers=none,
  basicstyle=\ttfamily\scriptsize\color{black},
}

% --- Estilos TikZ reutilizados ---
\tikzset{
  reg/.style={rectangle, draw=UFSBlue, fill=UFSBlue!8, minimum width=3.4cm,
    minimum height=0.44cm, font=\scriptsize\ttfamily, align=center},
  regmain/.style={reg, draw=UFSGray, fill=black!5},
  regbase/.style={reg, draw=UFSRed, fill=UFSRed!8},
  heap/.style={rectangle, draw=UFSRed, fill=UFSRed!8, rounded corners,
    font=\scriptsize\ttfamily, align=center, inner sep=4pt},
  zona/.style={rectangle, draw=UFSGray, minimum width=3.2cm,
    font=\scriptsize, align=center},
  rot/.style={font=\scriptsize\ttfamily, text=UFSGray},
  leg/.style={font=\scriptsize, text=UFSGray},
  seta/.style={->, thick, UFSBlue},
  setar/.style={->, thick, UFSRed},
  setag/.style={->, thick, UFSGray},
  caixa/.style={rectangle, draw=UFSGray, rounded corners, fill=black!3,
    font=\scriptsize, align=center, inner sep=4pt},
}

% Referencia da fonte no rodape do frame
\newcommand{\fonte}[1]{\vfill\hfill{\scriptsize\textcolor{UFSGray}{Fonte: #1}}}

% --- Metadados ---
\title{Recursividade e memória}
\subtitle{O que a recursão gasta, onde, e como não estourar}
\author{Prof.\ André Yoshiaki Kashiwabara}
\institute{DCOMP -- Universidade Federal de Sergipe\\
\texttt{andreyoshiaki@dcomp.ufs.br}}
\date{}

\begin{document}

\begin{frame}[plain]
  \titlepage
\end{frame}

\begin{frame}{O que você vai aprender hoje}
  \begin{center}
  \begin{tikzpicture}[
    item/.style={rectangle, draw=UFSBlue, fill=UFSBlue!10, rounded corners,
      minimum width=10.4cm, minimum height=0.6cm, font=\small}
  ]
    \node[item] (t1) at (0, 0)    {Até onde a pilha aguenta: o limite e o estouro};
    \node[item] (t2) at (0,-0.8)  {O que pesa em cada registro --- e o que morre no \texttt{return}};
    \node[item] (t3) at (0,-1.6)  {Quando nada fica para a volta: recursão de cauda};
    \node[item] (t4) at (0,-2.4)  {Uma tabela no heap para todas as chamadas};
    \node[item] (t5) at (0,-3.2)  {\texttt{malloc} em cada chamada: quem libera, e quando};
    \node[item] (t6) at (0,-4.0)  {Exercícios};
    \draw[seta] (t1) -- (t2);
    \draw[seta] (t2) -- (t3);
    \draw[seta] (t3) -- (t4);
    \draw[seta] (t4) -- (t5);
    \draw[seta] (t5) -- (t6);
  \end{tikzpicture}
  \end{center}
\end{frame}

% ==================================================================
\section{A pergunta que ficou}
% ==================================================================

\begin{frame}{O problema: a recursão não é de graça}
  \begin{columns}[T]
    \begin{column}{0.5\textwidth}
      \textbf{Da aula passada}\\[0.3em]
      \small
      Cada chamada ativa ocupa um \textbf{registro de ativação} na pilha, e
      nenhum é devolvido antes de a descida chegar ao caso base.

      \medskip
      A pergunta que ficou: \textcolor{UFSRed}{\texttt{soma(v, 10000000)}
      funciona?}
    \end{column}
    \begin{column}{0.5\textwidth}
      \textbf{O que vamos descobrir}\\[0.3em]
      \small
      A pilha tem um tamanho fixo --- e pequeno. Saber quanto cada chamada
      gasta, e o que mandar para o \textbf{heap}, é o que separa uma função
      recursiva que funciona de uma que derruba o programa.
    \end{column}
  \end{columns}
  \begin{block}{A pergunta de hoje}
    Quanta memória uma função recursiva gasta, em que lugar, e o que fazer
    quando ela não cabe?
  \end{block}
\end{frame}

\begin{frame}[fragile]{O programa de hoje: \texttt{memoria.c} (1/3)}
\begin{lstlisting}[style=codC]
/* memoria.c --- quanto de memoria a recursao gasta, e onde */
#include <stdio.h>
#include <stdlib.h>
/* A: a soma da aula passada: a adicao fica para a volta */
long soma(const int *v, int n)
{
    if (n == 0)
        return 0;
    return soma(v, n - 1) + v[n - 1];
}

/* B: a mesma soma, levando o parcial na ida: nada fica para a volta */
long soma_acum(const int *v, int n, long parcial)
{
    if (n == 0)
        return parcial;
    return soma_acum(v, n - 1, parcial + v[n - 1]);
}
\end{lstlisting}
\end{frame}

\begin{frame}[fragile]{\texttt{memoria.c} (2/3): uma tabela no heap}
\begin{lstlisting}[style=codC]
/* C: Fibonacci que anota cada resultado em uma tabela no heap */
long fib_memo(int n, long *memo)
{
    if (n < 2)
        return n;
    if (memo[n] == 0)
        memo[n] = fib_memo(n - 1, memo) + fib_memo(n - 2, memo);
    return memo[n];
}
\end{lstlisting}
  \begin{block}{Lembrete da aula passada}
    \small A \texttt{fib} ingênua recalculava os mesmos valores do zero:
    \texttt{fib(40)} fez 331\,160\,281 chamadas. Aqui, cada \texttt{fib(k)}
    calculado fica guardado em \texttt{memo[k]}.
  \end{block}
\end{frame}

\begin{frame}[fragile]{\texttt{memoria.c} (3/3): o \texttt{main}}
\begin{lstlisting}[style=codC, basicstyle=\ttfamily\fontsize{7}{7.9}\selectfont]
int main(int argc, char *argv[])
{
    int n = (argc > 1) ? atoi(argv[1]) : 1000;
    int *v = malloc(n * sizeof *v);         /* o vetor mora no heap */
    long *memo = calloc(91, sizeof *memo);  /* tabela zerada */
    if (v == NULL || memo == NULL) {
        free(v);
        free(memo);
        return 1;
    }
    for (int i = 0; i < n; i++)
        v[i] = 1;
    printf("C: fib(90) = %ld\n", fib_memo(90, memo));
    free(memo);

    printf("B: soma_acum(v, %d) = %ld\n", n, soma_acum(v, n, 0));
    printf("A: soma(v, %d) = %ld\n", n, soma(v, n));
    free(v);
    return 0;
}
\end{lstlisting}
\end{frame}

\begin{frame}{O que este programa imprime?}
  \begin{block}{Antes de compilar, escreva sua resposta}
    \begin{enumerate}
      \item Com \texttt{./memoria 100000}: quais linhas aparecem, e com que
        valores?
      \item E com \texttt{./memoria 1000000}? E com \texttt{10000000}?
      \item Se alguma execução falhar, qual linha sai por último?
      \item A linha \textbf{C} demora? A \texttt{fib} ingênua levaria quanto
        tempo para \texttt{fib(90)}?
      \item O vetor de 10 milhões de \texttt{int} ocupa 40\,MB. Ele é o
        problema?
    \end{enumerate}
  \end{block}
\end{frame}

\begin{frame}[fragile]{Compile, execute e compare}
\begin{lstlisting}[style=shellStyle]
$ gcc -Wall -Wextra memoria.c -o memoria
$ ./memoria 100000
C: fib(90) = 2880067194370816120
B: soma_acum(v, 100000) = 100000
A: soma(v, 100000) = 100000
$ ./memoria 1000000
C: fib(90) = 2880067194370816120
Segmentation fault: 11
$ echo $?
139
\end{lstlisting}
  \begin{center}
    \scriptsize\textcolor{UFSGray}{No Linux a mensagem é
    \texttt{Segmentation fault (core dumped)}. Com \texttt{10000000}, o mesmo
    resultado.}
  \end{center}
\end{frame}

% ==================================================================
\section{Até onde a pilha aguenta}
% ==================================================================

\begin{frame}{Por que 100 mil passou e 1 milhão não?}
  \begin{columns}[c]
    \begin{column}{0.48\textwidth}
      \small
      \begin{itemize}
        \item a linha \textbf{C} saiu: o heap deu conta da tabela;
        \item a linha \textbf{B} \emph{não} saiu: o programa morreu
          \emph{dentro} de \texttt{soma\_acum};
        \item o vetor de 1 milhão de \texttt{int} (4\,MB) foi alocado sem
          problema.
      \end{itemize}
    \end{column}
    \begin{column}{0.48\textwidth}
      \begin{block}{Pense antes de passar o slide}
        \small
        O que \texttt{soma\_acum(v, 1000000, 0)} empilha que
        \texttt{soma\_acum(v, 100000, 0)} não empilha?

        \medskip
        Existe um $n$ exato a partir do qual o programa cai? Do que ele
        depende?
      \end{block}
    \end{column}
  \end{columns}
\end{frame}

\begin{frame}{Definição: limite da pilha e estouro de pilha}
  \begin{block}{Limite da pilha}
    \small Cada programa recebe, ao iniciar, uma pilha de tamanho
    \textbf{fixo}, definido pelo sistema operacional. Em Linux e macOS o
    padrão é de cerca de \textbf{8\,MB}; o comando \texttt{ulimit -s} mostra
    o valor em KB.
  \end{block}
  \begin{block}{Estouro de pilha (\emph{stack overflow})}
    \small Quando os registros ativos passam desse limite, o próximo
    registro cai em memória que não pertence à pilha. O sistema
    interrompe o programa (\texttt{SIGSEGV}, \emph{segmentation fault}).
  \end{block}
  \begin{exampleblock}{Intuição}
    \small profundidade máxima $\approx$ tamanho da pilha $\div$ tamanho de
    um registro
  \end{exampleblock}
  \fonte{Kerrisk (2010), §36.3~\cite{kerrisk2010linux};
    Bryant \& O'Hallaron (2015), §3.7.1~\cite{bryant2015computer}}
\end{frame}

\begin{frame}{A conta: quantos registros cabem em 8\,MB}
  \begin{columns}[c]
    \begin{column}{0.5\textwidth}
      \begin{center}
        \small
        \begin{tabular}{rrc}
          \toprule
          $n$ & pilha necessária & cabe?\\
          \midrule
          1\,000      & 48\,KB  & sim\\
          100\,000    & 4,8\,MB & sim\\
          1\,000\,000 & 48\,MB  & \textcolor{UFSRed}{não}\\
          10\,000\,000 & 480\,MB & \textcolor{UFSRed}{não}\\
          \bottomrule
        \end{tabular}

        \medskip
        \scriptsize registros de 48 bytes, compilado sem otimização
      \end{center}
    \end{column}
    \begin{column}{0.46\textwidth}
      \small
      Nesta máquina, \texttt{ulimit -s} dá 8\,176\,KB
      = 8\,372\,224 bytes.

      \medskip
      Por busca binária, o maior $n$ que passa fica em torno de \textbf{174 mil} (varia um pouco a cada execução):
      \[
        8\,372\,224 \div 174\,000 \approx 48 \text{ bytes por registro.}
      \]
      O limite não é um número mágico: é uma divisão.
    \end{column}
  \end{columns}
  \fonte{Kerrisk (2010), §36.3~\cite{kerrisk2010linux}; medido com
    \texttt{memoria.c}}
\end{frame}

\begin{frame}{E o vetor de 40\,MB? Por que ele não estourou nada?}
  \begin{block}{Pense antes de passar o slide}
    \small No \texttt{main}, \texttt{v} foi criado com \texttt{malloc}. O que
    aconteceria se ele fosse declarado assim?
    \begin{center}
      \texttt{int v[10000000];}
    \end{center}
    E por que 40\,MB no heap não são problema, se 48\,MB na pilha são?
  \end{block}
\end{frame}

\begin{frame}{A resposta: pilha e heap são regiões diferentes}
  \begin{columns}[c]
    \begin{column}{0.44\textwidth}
      \begin{center}
      \begin{tikzpicture}
        \node[zona, minimum height=1.0cm, fill=UFSBlue!8, draw=UFSBlue]
          (p) at (0,0) {\textbf{pilha}\\registros: \texttt{n}, \texttt{parcial}};
        \node[zona, minimum height=0.9cm, fill=white, dashed]
          (l) at (0,-0.95) {\textcolor{UFSGray}{livre}};
        \node[zona, minimum height=1.0cm, fill=UFSRed!8, draw=UFSRed]
          (h) at (0,-1.9) {\textbf{heap}\\\texttt{v}, \texttt{memo}};
        \node[zona, minimum height=0.5cm, fill=black!5] (d) at (0,-2.65)
          {dados estáticos};
        \node[zona, minimum height=0.5cm, fill=black!5] (c) at (0,-3.15)
          {código};
        \draw[seta] (p.south east) ++(0.15,0.3) -- ++(0,-0.6);
        \draw[setar] (h.north east) ++(0.15,-0.3) -- ++(0,0.6);
        \node[leg, anchor=west] at (1.75,0.1) {limite fixo};
        \node[leg, anchor=west] at (1.75,-1.8) {cresce sob};
        \node[leg, anchor=west] at (1.75,-2.1) {demanda};
        \node[leg, anchor=east] at (-1.7,0.4) {endereços};
        \node[leg, anchor=east] at (-1.7,0.1) {altos};
      \end{tikzpicture}
      \end{center}
    \end{column}
    \begin{column}{0.54\textwidth}
      \small
      \texttt{int v[10000000]} como local vai para o registro do
      \texttt{main}: 40\,MB numa pilha de 8\,MB. O programa cai
      \textbf{antes} da primeira chamada.

      \medskip
      O heap cresce sob demanda, até o limite da memória do processo.
      Se não houver espaço, \texttt{malloc} devolve \texttt{NULL}
      --- um erro que dá para \emph{testar}. O estouro de pilha não avisa.
    \end{column}
  \end{columns}
  \fonte{Kerrisk (2010), §6.3~\cite{kerrisk2010linux};
    Bryant \& O'Hallaron (2015), §9.9~\cite{bryant2015computer}}
\end{frame}

\begin{frame}{Recapitulando: até onde a pilha aguenta}
  \begin{block}{O que aprendemos}
    \begin{itemize}
      \item a pilha tem tamanho fixo, da ordem de 8\,MB;
      \item a profundidade máxima é o limite dividido pelo tamanho de um
        registro;
      \item passar do limite é um \emph{segmentation fault}, sem aviso e
        sem chance de recuperação;
      \item dados grandes vão para o heap, onde a falta de espaço vira um
        \texttt{NULL} testável.
    \end{itemize}
  \end{block}
\end{frame}

% ==================================================================
\section{O que pesa em cada registro}
% ==================================================================

\begin{frame}[fragile]{Um vetor local em cada nível}
  \begin{columns}[c]
    \begin{column}{0.52\textwidth}
\begin{lstlisting}[style=tinyC]
/* gorda.c */
int desce(int n)
{
    char linha[4096];   /* 4 KB */
    linha[0] = (char) n;
    if (n == 0)
        return linha[0];
    return desce(n - 1) + linha[0];
}
\end{lstlisting}
    \end{column}
    \begin{column}{0.44\textwidth}
      \begin{block}{Pense antes de passar o slide}
        \small \texttt{soma\_acum} aguentou cerca de 174 mil níveis. Quantos
        níveis \texttt{desce} aguenta?
        \begin{itemize}
          \item uns 170 mil?
          \item uns 20 mil?
          \item uns 2 mil?
        \end{itemize}
      \end{block}
    \end{column}
  \end{columns}
\end{frame}

\begin{frame}{A resposta: as locais se multiplicam pela profundidade}
  \begin{columns}[c]
    \begin{column}{0.5\textwidth}
      \small
      Medido: \texttt{desce} cai a partir de \textbf{2\,010} níveis.
      \[
        8\,372\,224 \div 2\,010 \approx 4\,165 \text{ bytes}
      \]
      = os 4\,096 de \texttt{linha} + o resto do registro.

      \medskip
      Um vetor local existe uma vez \emph{por chamada ativa}. Em uma função
      recursiva, ele é pago a cada nível.
    \end{column}
    \begin{column}{0.46\textwidth}
      \begin{block}{O que fazer}
        \small
        \begin{itemize}
          \item alocar o espaço \textbf{uma vez}, no heap, fora da
            recursão;
          \item passar o ponteiro para baixo como parâmetro;
          \item cada registro paga só 8 bytes do ponteiro.
        \end{itemize}
      \end{block}
    \end{column}
  \end{columns}
  \fonte{Bryant \& O'Hallaron (2015), §3.7.4~\cite{bryant2015computer};
    medido com \texttt{gorda.c}}
\end{frame}

\begin{frame}[fragile]{E se a função devolver o endereço de uma local?}
  \begin{columns}[c]
    \begin{column}{0.52\textwidth}
\begin{lstlisting}[style=tinyC]
char *rotulo(int n)
{
    char texto[16];
    snprintf(texto, sizeof texto,
             "nivel %d", n);
    return texto;
}

int main(void)
{
    char *a = rotulo(1);
    char *b = rotulo(2);
    printf("%s | %s\n", a, b);
    return 0;
}
\end{lstlisting}
    \end{column}
    \begin{column}{0.44\textwidth}
      \begin{block}{Pense antes de passar o slide}
        \small O que o \texttt{printf} imprime?

        \medskip
        Onde está \texttt{texto} depois que \texttt{rotulo} retorna?
      \end{block}
    \end{column}
  \end{columns}
\end{frame}

\begin{frame}[fragile]{A resposta: o registro morreu com o \texttt{return}}
\begin{lstlisting}[style=shellStyle]
$ gcc -Wall -Wextra rotulo.c -o rotulo
rotulo.c:6:12: warning: address of stack memory associated with local
      variable 'texto' returned [-Wreturn-stack-address]
$ ./rotulo
??? | ???        <- bytes ilegiveis, outros a cada execucao
\end{lstlisting}
  \begin{block}{}
    \small A vida de uma local termina quando a chamada termina. O ponteiro
    continua apontando para onde ela estava --- uma área que a próxima
    chamada reaproveita. O resultado é imprevisível; o compilador avisa
    (no gcc: \texttt{-Wreturn-local-addr}).
  \end{block}
  \begin{center}
    \small Para devolver texto: receba o espaço do chamador, ou devolva
    memória do \texttt{malloc} --- e documente quem dá \texttt{free}.
  \end{center}
  \fonte{ISO/IEC 9899:2018, §6.2.4~\cite{iso2018c17};
    Bryant \& O'Hallaron (2015), §9.11~\cite{bryant2015computer}}
\end{frame}

\begin{frame}{Recapitulando: o que pesa em cada registro}
  \begin{block}{O que aprendemos}
    \begin{itemize}
      \item parâmetros e locais --- inclusive vetores --- moram no
        registro;
      \item numa função recursiva, cada local é paga uma vez por nível;
      \item espaço grande: aloque uma vez no heap e passe o ponteiro;
      \item o registro deixa de existir no \texttt{return}: nunca devolva
        o endereço de uma local.
    \end{itemize}
  \end{block}
\end{frame}

% ==================================================================
\section{Recursão de cauda}
% ==================================================================

\begin{frame}[fragile]{Por que, com \texttt{-O2}, a linha \textbf{B} passou?}
\begin{lstlisting}[style=shellStyle]
$ gcc -O2 -Wall -Wextra memoria.c -o memoria
$ ./memoria 1000000
C: fib(90) = 2880067194370816120
B: soma_acum(v, 1000000) = 1000000
Segmentation fault: 11
\end{lstlisting}
  \begin{block}{Pense antes de passar o slide}
    \small Mesmo programa, mesmo $n$. Só mudou a opção de otimização.

    \medskip
    Por que \texttt{soma\_acum} agora aguenta um milhão de níveis e
    \texttt{soma} não? O que uma tem que a outra não tem?
  \end{block}
\end{frame}

\begin{frame}{Definição: chamada de cauda}
  \begin{block}{Chamada de cauda}
    \small Uma chamada é \textbf{de cauda} quando é a última coisa que a
    função faz: o valor que ela devolve é devolvido direto, sem nenhuma
    conta depois. Uma função recursiva cujas chamadas são todas de cauda é
    uma \textbf{recursão de cauda}.
  \end{block}
  \begin{exampleblock}{Intuição}
    \small Se nada fica pendente, o registro de quem chamou não serve para
    mais nada: o compilador pode \emph{reaproveitá-lo} e trocar a chamada
    por um salto --- a recursão vira um laço.
  \end{exampleblock}
  \begin{center}
    \small
    \texttt{return soma(v, n - 1) + v[n - 1];} \quad
    \textcolor{UFSRed}{a soma fica para depois}\\[0.2em]
    \texttt{return soma\_acum(v, n - 1, parcial + v[n - 1]);} \quad
    \textcolor{UFSBlue}{nada fica}
  \end{center}
  \fonte{Abelson et al.\ (1996), §1.2.1~\cite{abelson1996structure}}
\end{frame}

\begin{frame}{O que fica pendente em cada registro}
  \begin{columns}[T]
    \begin{column}{0.48\textwidth}
      \begin{center}
      \textcolor{UFSRed}{\textbf{\texttt{soma(v, 3)}}}\\[0.4em]
      \begin{tikzpicture}
        \node[regmain] (m)  at (0,0)    {main};
        \node[reg]     (a)  at (0,0.45) {n = 3 \ espera ? + v[2]};
        \node[reg]     (b)  at (0,0.90) {n = 2 \ espera ? + v[1]};
        \node[reg]     (c)  at (0,1.35) {n = 1 \ espera ? + v[0]};
        \node[regbase] (d)  at (0,1.80) {n = 0 \ devolve 0};
      \end{tikzpicture}

      \scriptsize cada registro guarda uma conta por terminar:\\
      todos precisam continuar vivos
      \end{center}
    \end{column}
    \begin{column}{0.48\textwidth}
      \begin{center}
      \textcolor{UFSBlue}{\textbf{\texttt{soma\_acum(v, 3, 0)}}}\\[0.4em]
      \begin{tikzpicture}
        \node[regmain] (m)  at (0,0)    {main};
        \node[reg]     (a)  at (0,0.45) {n \ \ parcial};
      \end{tikzpicture}

      \vspace{0.3em}
      {\scriptsize\ttfamily
      \begin{tabular}{rrl}
        n & parcial & \\
        \midrule
        3 & 0 & \\
        2 & 1 & \\
        1 & 2 & \\
        0 & 3 & \textcolor{UFSRed}{devolve 3}
      \end{tabular}}

      \scriptsize o resultado parcial desce junto:\\
      um registro só, atualizado no lugar
      \end{center}
    \end{column}
  \end{columns}
  \fonte{Abelson et al.\ (1996), §1.2.1~\cite{abelson1996structure}}
\end{frame}

\begin{frame}[fragile]{O que o compilador fez com \texttt{soma\_acum}}
  \begin{columns}[T]
    \begin{column}{0.48\textwidth}
      \textcolor{UFSBlue}{\textbf{O que você escreveu}}
\begin{lstlisting}[style=tinyC]
long soma_acum(const int *v, int n,
               long parcial)
{
    if (n == 0)
        return parcial;
    return soma_acum(v, n - 1,
                     parcial + v[n - 1]);
}
\end{lstlisting}
    \end{column}
    \begin{column}{0.48\textwidth}
      \textcolor{UFSBlue}{\textbf{O que o \texttt{-O2} gerou (em C)}}
\begin{lstlisting}[style=tinyC]
long soma_acum(const int *v, int n,
               long parcial)
{
    while (n != 0) {
        parcial = parcial + v[n - 1];
        n = n - 1;
    }
    return parcial;
}
\end{lstlisting}
    \end{column}
  \end{columns}
  \begin{block}{}
    \small Os parâmetros viram as variáveis do laço; a chamada vira
    ``volte ao início''. Um registro só, qualquer que seja $n$.
  \end{block}
  \fonte{Abelson et al.\ (1996), §1.2.1~\cite{abelson1996structure};
    Bryant \& O'Hallaron (2015), §3.7~\cite{bryant2015computer}}
\end{frame}

\begin{frame}{Não conte com isso}
  \begin{alertblock}{A otimização não é garantida}
    \small
    \begin{itemize}
      \item o padrão C não exige que chamadas de cauda virem laço;
      \item sem \texttt{-O2} --- como na depuração --- \texttt{soma\_acum}
        caiu em cerca de 174 mil níveis;
      \item outro compilador ou outra versão pode fazer diferente, para mais
        ou para menos;
      \item uma mudança pequena (um \texttt{printf} depois da chamada) tira a
        chamada da cauda.
    \end{itemize}
  \end{alertblock}
  \begin{block}{Regra prática}
    \small Se a profundidade pode chegar a dezenas de milhares, escreva o
    laço você mesmo. A versão de cauda é o meio do caminho: já tem a cara
    do laço.
  \end{block}
  \fonte{ISO/IEC 9899:2018, §6.5.2.2~\cite{iso2018c17}; medido com
    \texttt{memoria.c} (clang, macOS)}
\end{frame}

\begin{frame}{Recapitulando: recursão de cauda}
  \begin{block}{O que aprendemos}
    \begin{itemize}
      \item numa chamada de cauda nada fica pendente no registro de quem
        chamou;
      \item o resultado parcial desce como parâmetro (um acumulador);
      \item com otimização, o compilador \emph{pode} transformar a
        recursão em laço;
      \item ``pode'' não é ``vai'': para profundidades grandes, use o laço.
    \end{itemize}
  \end{block}
\end{frame}

% ==================================================================
\section{Uma tabela para todas as chamadas}
% ==================================================================

\begin{frame}{Por que \texttt{fib(90)} saiu na hora?}
  \begin{columns}[c]
    \begin{column}{0.5\textwidth}
      \small
      A \texttt{fib} ingênua faz $2F(n+1) - 1$ chamadas. Para $n = 90$:
      \[
        \approx 9{,}3 \times 10^{18} \text{ chamadas}
      \]
      A um bilhão de chamadas por segundo, quase \textbf{300 anos}.
    \end{column}
    \begin{column}{0.46\textwidth}
      \begin{block}{Pense antes de passar o slide}
        \small Quantas chamadas \texttt{fib\_memo(90, memo)} faz?

        \medskip
        Quantas tabelas \texttt{memo} existem durante a execução --- uma
        por chamada, ou uma só?
      \end{block}
    \end{column}
  \end{columns}
  \fonte{Deitel \& Deitel (2011), §5.15~\cite{deitel2011como}}
\end{frame}

\begin{frame}{Definição: memoização}
  \begin{block}{Memoização}
    \small Guardar o resultado de cada chamada já calculada e, na próxima vez
    que ela for pedida, \textbf{consultar} em vez de recalcular.
  \end{block}
  \begin{columns}[c]
    \begin{column}{0.56\textwidth}
      \begin{center}
      \begin{tikzpicture}
        \node[regmain] (m)  at (0,0)    {main \ memo = \ldots};
        \node[reg]     (a)  at (0,0.45) {fib\_memo \ n = 90};
        \node[reg]     (b)  at (0,0.90) {fib\_memo \ n = 89};
        \node[reg]     (c)  at (0,1.35) {\ldots};
        \node[heap]    (t)  at (3.3,0.7) {memo[0..90]\\91 \texttt{long}\\728 bytes};
        \draw[setar] (m.east) -- (t);
        \draw[setar] (a.east) -- (t);
        \draw[setar] (b.east) -- (t);
        \node[leg] at (3.3,-0.1) {heap: uma só};
      \end{tikzpicture}
      \end{center}
    \end{column}
    \begin{column}{0.4\textwidth}
      \small Cada registro tem só uma \emph{cópia do ponteiro}; a tabela é
      uma só. Resultado: \textbf{179} chamadas ($2n - 1$) em vez de
      $9{,}3 \times 10^{18}$.
    \end{column}
  \end{columns}
  \fonte{Michie (1968)~\cite{michie1968memo};
    Cormen et al.\ (2009), §15.3~\cite{cormen2009introduction}}
\end{frame}

\begin{frame}{Quem aloca, quem libera}
  \begin{columns}[T]
    \begin{column}{0.48\textwidth}
      \begin{block}{Fora da recursão}
        \small A tabela precisa existir \textbf{antes} da primeira chamada e
        \textbf{depois} da última. Por isso o \texttt{main} aloca, passa o
        ponteiro e libera no fim.
      \end{block}
    \end{column}
    \begin{column}{0.48\textwidth}
      \begin{block}{Por que \texttt{calloc}}
        \small O zero significa ``ainda não calculado''. Com
        \texttt{malloc}, a tabela viria com lixo --- e o lixo seria lido como
        resultado pronto.
      \end{block}
    \end{column}
  \end{columns}
  \medskip
  \begin{center}
    \small
    \begin{tabular}{@{}lll@{}}
      \toprule
      é de\ldots & vai para & exemplo\\
      \midrule
      uma chamada só & a pilha (local) & \texttt{n}, \texttt{parcial}\\
      todas as chamadas & o heap, alocado fora & \texttt{memo}, \texttt{v}\\
      \bottomrule
    \end{tabular}
  \end{center}
  \fonte{Kernighan \& Ritchie (1989), §7.8.5~\cite{kernighan1989c};
    Cormen et al.\ (2009), §15.3~\cite{cormen2009introduction}}
\end{frame}

\begin{frame}{Recapitulando: uma tabela para todas as chamadas}
  \begin{block}{O que aprendemos}
    \begin{itemize}
      \item memoização troca memória por tempo: 728 bytes no lugar de
        séculos;
      \item o que é compartilhado por todas as chamadas vai para o heap,
        alocado fora da recursão;
      \item cada registro carrega só o ponteiro, não a tabela;
      \item \texttt{calloc} dá uma tabela zerada --- e o zero precisa ter
        significado.
    \end{itemize}
  \end{block}
\end{frame}

% ==================================================================
\section{\texttt{malloc} em cada chamada}
% ==================================================================

\begin{frame}[fragile]{Uma função recursiva que devolve memória do heap}
  \begin{columns}[c]
    \begin{column}{0.6\textwidth}
\begin{lstlisting}[style=codC]
/* invertida.c: copia de s de tras para frente */
char *invertida(const char *s)
{
    if (*s == '\0') {              /* caso base */
        char *r = aloca(1);
        r[0] = '\0';
        return r;
    }
    char *resto = invertida(s + 1);
    size_t k = strlen(resto);
    char *r = aloca(k + 2);
    strcpy(r, resto);   /* o resto invertido... */
    r[k] = *s;          /* ...e a minha letra   */
    r[k + 1] = '\0';
    libera(resto);      /* na volta */
    return r;
}
\end{lstlisting}
    \end{column}
    \begin{column}{0.36\textwidth}
      \small
      \texttt{aloca} e \texttt{libera} chamam \texttt{malloc} e
      \texttt{free}, e contam quantos bytes estão no heap agora e no
      máximo (o \emph{pico}).

      \medskip
      Cada nível recebe um bloco da chamada de baixo, monta um bloco maior e
      devolve.
    \end{column}
  \end{columns}
\end{frame}

\begin{frame}{Quanto heap isso usa?}
  \begin{block}{Antes de executar, escreva sua resposta}
    \begin{enumerate}
      \item Para uma string de $n = 10\,000$ letras, quantos bytes ficam no
        heap no \textbf{pico}? Uns 20 mil? Uns 10 milhões?
      \item E se apagarmos a linha \texttt{libera(resto);}? Quanto fica no
        heap no fim?
      \item Sem o \texttt{libera}, o resultado impresso muda?
    \end{enumerate}
  \end{block}
\end{frame}

\begin{frame}[fragile]{Execute e compare}
  \begin{center}
    \small
    \begin{tabular}{@{}rrrr@{}}
      \toprule
      & \multicolumn{1}{c}{com \texttt{libera(resto)}}
      & \multicolumn{2}{c}{sem \texttt{libera(resto)}}\\
      \cmidrule(lr){2-2}\cmidrule(lr){3-4}
      $n$ & pico & pico & perdido no fim\\
      \midrule
      10      & 21       & 66           & 55\\
      1\,000  & 2\,001   & 501\,501     & 500\,500\\
      10\,000 & 20\,001  & 50\,015\,001 & 50\,005\,000\\
      \bottomrule
    \end{tabular}
  \end{center}
  \begin{block}{}
    \small Sem o \texttt{libera}, o resultado impresso é \textbf{o mesmo} ---
    nada avisa. Para inverter 10\,KB de texto, o programa deixa 50\,MB
    perdidos no heap.
  \end{block}
  \begin{center}
    \scriptsize\textcolor{UFSGray}{``perdido'' = pico menos o bloco do
    resultado, que o \texttt{main} ainda usa e libera.}
  \end{center}
  \fonte{medido com \texttt{invertida.c}}
\end{frame}

\begin{frame}{Por que $2n$ com o \texttt{free} e $n^2/2$ sem ele}
  \begin{columns}[c]
    \begin{column}{0.5\textwidth}
      \begin{center}
      \begin{tikzpicture}[x=0.45cm, y=0.32cm]
        \foreach \k in {1,...,6} {
          \fill[UFSRed!20, draw=UFSRed] (\k-1,0) rectangle (\k,\k);
        }
        \draw[UFSBlue, very thick] (4,0) rectangle (6,6);
        \node[leg, below] at (3,0) {blocos de 1, 2, \ldots, $n+1$ bytes};
        \node[leg, text=UFSBlue, align=left, anchor=west] at (6.3,4.5)
          {com \texttt{free}:\\só os 2 últimos\\vivos juntos};
        \node[leg, text=UFSRed, align=left, anchor=west] at (6.3,1.2)
          {sem \texttt{free}:\\todos ficam};
      \end{tikzpicture}
      \end{center}
    \end{column}
    \begin{column}{0.46\textwidth}
      \small
      Com o \texttt{free}, enquanto um nível monta o seu bloco ($k+2$
      bytes), só o do nível de baixo ($k+1$) ainda existe: pico
      $\approx 2n$.

      \medskip
      Sem o \texttt{free}, nenhum bloco sai:
      \[
        1 + 2 + \cdots + (n+1) = \tfrac{(n+1)(n+2)}{2}
      \]
      É um \textbf{vazamento de memória} que cresce com o quadrado de $n$.
    \end{column}
  \end{columns}
  \fonte{King (2008), §17.4~\cite{king2008c};
    Bryant \& O'Hallaron (2015), §9.11~\cite{bryant2015computer}}
\end{frame}

\begin{frame}{A regra: quem recebe o bloco é dono dele}
  \begin{columns}[T]
    \begin{column}{0.48\textwidth}
      \begin{block}{O contrato}
        \small \texttt{invertida} devolve um bloco do heap. Quem chama
        \emph{recebe a posse}: precisa usá-lo e depois dar \texttt{free}
        --- seja o \texttt{main}, seja o nível de cima da própria recursão.
      \end{block}
    \end{column}
    \begin{column}{0.48\textwidth}
      \begin{block}{Ida e volta, de novo}
        \small O \texttt{libera(resto)} vem \textbf{depois} da chamada:
        roda na volta, quando o bloco já foi copiado. Antes do
        \texttt{strcpy}, seria um uso de memória já liberada.
      \end{block}
    \end{column}
  \end{columns}
  \medskip
  \begin{center}
    \small Escreva o contrato no comentário da função:\\[0.2em]
    \texttt{/* devolve bloco do heap; quem chama deve dar free */}
  \end{center}
  \fonte{Kernighan \& Ritchie (1989), §7.8.5~\cite{kernighan1989c};
    King (2008), §17.4~\cite{king2008c}}
\end{frame}

\begin{frame}{Recapitulando: \texttt{malloc} em cada chamada}
  \begin{block}{O que aprendemos}
    \begin{itemize}
      \item cada nível que aloca precisa liberar o que não vai devolver;
      \item o \texttt{free} do bloco recebido vem depois de usá-lo --- na
        volta;
      \item um vazamento dentro da recursão se multiplica pelo número de
        chamadas;
      \item o resultado certo não prova que a memória está certa: meça.
    \end{itemize}
  \end{block}
\end{frame}

% ==================================================================
\section{Mãos à obra}
% ==================================================================

\begin{frame}{Altere o programa}
  \small
  \begin{block}{Tarefa 1}
    Descubra o limite da \emph{sua} máquina: o maior $n$ que
    \texttt{./memoria} aceita sem e com \texttt{-O2}. Rode
    \texttt{ulimit -s} e calcule o tamanho de cada registro.
  \end{block}
  \begin{block}{Tarefa 2}
    Rode \texttt{ulimit -s 65536} e repita \texttt{./memoria 1000000}.
    Depois ponha um \texttt{int rascunho[32]} (usado) em
    \texttt{soma\_acum}. Preveja o novo limite antes de medir.
  \end{block}
  \begin{block}{Tarefa 3}
    Troque \texttt{calloc} por \texttt{malloc} na tabela. O que muda? E se
    a tabela tiver 90 posições em vez de 91?
  \end{block}
\end{frame}

\begin{frame}{Exercícios (1/2)}
  \footnotesize
  \begin{block}{Exercício 1 --- somar pelas metades}
    Escreva \texttt{long soma\_metades(const int *v, int ini, int fim)}
    que soma \texttt{v[ini..fim-1]} somando a metade esquerda e a direita.
    Ela funciona com 10 milhões de elementos sem \texttt{-O2}. Por quê?
  \end{block}
  \begin{block}{Exercício 2 --- o vetor auxiliar vai para o heap}
    Uma ordenação recursiva declara \texttt{int aux[fim - ini]} a cada
    chamada e cai com 10 milhões de elementos. Aloque \texttt{aux}
    \textbf{uma vez}, no heap, e passe-o para baixo.
  \end{block}
  \fonte{Feofiloff (2009), cap.~2~\cite{feofiloff2009algoritmos}}
\end{frame}

\begin{frame}{Exercícios (2/2)}
  \footnotesize
  \begin{block}{Exercício 3 --- máximo com acumulador}
    \texttt{maximo} em recursão de cauda, levando o maior até agora como
    parâmetro. Teste com 1 milhão de elementos, sem e com \texttt{-O2}.
  \end{block}
  \begin{block}{Exercício 4 --- combinações com tabela}
    $C(n,k) = C(n-1,k-1) + C(n-1,k)$. Calcule $C(60, 30)$ com uma tabela
    $(n+1) \times (k+1)$ em bloco único no heap (\texttt{memo[i * nc + j]}).
  \end{block}
  \begin{block}{Exercício 5 --- repetir sem vazar}
    \texttt{char *repete(const char *s, int k)} devolve \texttt{s} repetida
    \texttt{k} vezes, num bloco do heap. No fim, o heap volta a zero.
  \end{block}
  \fonte{Cormen et al.\ (2009), §15.3~\cite{cormen2009introduction};
    King (2008), §17.4~\cite{king2008c}}
\end{frame}

\begin{frame}{Desafio}
  \begin{block}{Pintar uma região: troque a pilha do sistema por uma sua}
    \small Uma grade de \texttt{l} $\times$ \texttt{c} caracteres, alocada
    no heap, tem células livres (\texttt{.}) e paredes (\texttt{X}).
    \begin{itemize}
      \item \texttt{pinta\_rec}: a partir de \texttt{(i, j)}, marca a célula
        com \texttt{\#} e chama a si mesma para os quatro vizinhos; devolve
        quantas pintou;
      \item \texttt{pinta\_pilha}: o mesmo, sem recursão, com uma pilha de
        posições no heap que cresce com \texttt{realloc}.
    \end{itemize}
    \begin{itemize}
      \item as duas concordam em grades pequenas;
      \item numa grade $3000 \times 3000$, a recursiva cai --- mesmo com
        \texttt{-O2}. Explique por quê;
      \item a versão com pilha própria pinta as 4\,500\,000 células e
        libera tudo.
    \end{itemize}
  \end{block}
  \fonte{Bryant \& O'Hallaron (2015), §9.9~\cite{bryant2015computer}}
\end{frame}

\begin{frame}{Fechamento}
  \begin{columns}[T]
    \begin{column}{0.5\textwidth}
      \textbf{O que mudou hoje}\\[0.3em]
      \small A pergunta da aula passada tem resposta: dez milhões de
      registros não cabem. Mas a soma de dez milhões cabe --- em um laço, ou
      em uma recursão que desce só $\log_2 n$ níveis.
    \end{column}
    \begin{column}{0.5\textwidth}
      \textbf{O preço}\\[0.3em]
      \small A pilha é pequena, e estourá-la não avisa. O heap é grande, mas
      cada bloco precisa de um dono que dê \texttt{free}.
    \end{column}
  \end{columns}
  \begin{block}{Para levar}
    \begin{itemize}
      \small
      \item profundidade máxima $\approx$ pilha $\div$ registro;
      \item vetores locais pesam em cada nível: aloque fora e passe o
        ponteiro;
      \item recursão de cauda \emph{pode} virar laço --- não conte com isso;
      \item o que é de todas as chamadas vai para o heap; quem recebe um
        bloco o libera.
    \end{itemize}
  \end{block}
\end{frame}

\section*{Referências}
\begin{frame}[allowframebreaks]{Referências}
  \nocite{*}
  \bibliographystyle{plain}
  \bibliography{referencias}
\end{frame}

\end{document}
